%%%%%%%%%%%%%%%%%%%%%%%%%%%%%%%%%%%%%%%%%%%%%%%%%%%%%%%%%%%%%%%%%%%%
% Einleitung
%%%%%%%%%%%%%%%%%%%%%%%%%%%%%%%%%%%%%%%%%%%%%%%%%%%%%%%%%%%%%%%%%%%%

\chapter{Einleitung}
\label{Einleitung}

%%%%%%%%%%%%%%%%%%%%%geschrieben von Luis%%%%%%%%%%%%%%%%%%%%%
%%%% nur exemplarisch, kann geändert und angepasst werden%%%%
Die vorliegende Studienarbeit befasst sich im ersten Schritt mit der fundierten Erarbeitung der verwendeten methodischen und technologischen Grundlagen. Darauf aufbauend erfolgt eine detaillierte Analyse der Ausgangslage, welche explizit an die Ergebnisse und den Stand der vorherigen Studienarbeit anknüpft. Bestandteil dieser Evaluation sind eine umfassende Ist- und Laufzeitanalyse sowie die Ausarbeitung daraus abgeleiteter Optimierungsvorschläge.
\\
\\Auf Basis dieser Erkenntnisse widmet sich die Arbeit der Konzeptionierung neuer Lösungsansätze. Ein zentraler Fokus liegt hierbei auf der Optimierung und Erweiterung der bestehenden Systemarchitektur. Dies beinhaltet insbesondere die effiziente Speicherung und Instanziierung von Punktwolken. Zur erweiterten visuellen Auswertung wird eine Funktion zur wahlweisen einfachen oder doppelten Anzeige der Punktwolken integriert, wobei der Wechsel zwischen diesen beiden Darstellungsmodi nutzerfreundlich über eine neu implementierte Auswahltaste erfolgt.
\\
\\Ein weiterer inhaltlicher Schwerpunkt ist die informationstechnische Umsetzung einer Befehlspipeline sowie relevanter Safety-Features. In diesem Zuge werden die zugrunde liegende Softwarestruktur und die angewandten Design-Patterns detailliert beleuchtet. Nach einer kritischen Abwägung und Einordnung verschiedener Ansätze zur Safety-Implementierung wird die Wahl der finalen Methode sachlich begründet und deren praktische Umsetzung dargelegt. Den Abschluss der Arbeit bildet ein Ausblick, der potenzielle Erweiterungen, gezielte Verbesserungsvorschläge und konkrete Anknüpfungspunkte für nachfolgende Studienarbeiten aufzeigt.

